\bookchapter{บทที่ 27: พ่อแม่ก็ใช้ง่าย—คู่มือปรับแต่งมือถือสำหรับผู้สูงอายุอย่างครบถ้วน}

\chapterintro{พูดตรงๆ ตอนเขียนบทนี้ ฉันนึกถึงพ่อกับแม่ของฉัน ทุกครั้งที่มือถือของแม่ฉันเด้งข้อความ “อัปเดตระบบ” ขึ้นมา แม่ไม่กล้ากด เพราะกลัวว่าจะทำให้มือถือพัง ส่วนพ่อของฉันตั้งค่าตัวอักษร WeChat เล็กกว่ามดอีก แต่ก็ไม่รู้วิธีขยายให้ใหญ่ขึ้น ต้องมัวเพ่งตาอ่านอยู่แบบนั้นตั้ง 2 ปี นี่ไม่ใช่เพราะพวกท่าน “โง่” แต่เป็นเพราะมือถือตั้งค่าเริ่มต้นโดยถือว่าทุกคนเป็นคนหนุ่มสาววัยยี่สิบกว่าๆ ทั้งหมด ในบทนี้ ฉันจะใช้คำพูดที่ตรงไปตรงมาที่สุด เพื่อไขปัญหาที่ทำให้ผู้สูงอายุปวดหัวที่สุดในการใช้มือถือทีละข้อ}

\booksection{ตั้งใจให้ถูกต้องก่อน: นี่ไม่ใช่เรื่อง “ผู้สูงอายุใช้ไม่เป็น” แต่เป็นเรื่อง “มือถือถูกออกแบบมาไม่ดี”}

ก่อนจะตั้งค่าใดๆ มีแนวคิดหนึ่งที่ต้องอธิบายให้ชัดเจนก่อน:\textbf{อินเทอร์เฟซของสมาร์ทโฟปัจจุบัน ตั้งแต่เริ่มออกแบบก็ไม่ได้คำนึงถึงผู้ใช้สูงอายุมากนัก}หน้าต่างป๊อปอัปเล็กจนมองไม่เห็นชัด ท่าทางการสัมผัสต้องการความแม่นยำสูง ไอคอนหน้าตาคล้ายกันไปหมด—ปัญหาเหล่านี้ไม่ใช่ความผิดของผู้ใช้ แต่เป็นปัญหาของผลิตภัณฑ์ ดังนั้นการพาผู้ใหญ่ใช้มือถือ สิ่งแรกไม่ใช่การสอน แต่คือการช่วยพวกเขา\textbf{ปรับสภาพแวดล้อมให้พร้อม}เมื่อสภาพแวดล้อมถูกปรับให้เหมาะสมแล้ว ปัญหาก็จะหายไปเกินครึ่ง

การพาผู้ใหญ่ใช้มือถือมีกฎเหล็ก 3 ข้อ: (1)\textbf{อย่าเร่งเร้า}—คุณใช้สมาร์ทโฟนมาสิบปี แต่พวกเขาอาจจะใช้มาแค่สามเดือน (2)\textbf{ใช้การเปรียบเทียบ}—“Wi-Fi เหมือนกับท่อน้ำในบ้าน เปิดก๊อกน้ำถึงจะมีน้ำไหล” การเปรียบเทียบแบบนี้เข้าใจง่ายกว่าคำศัพท์เทคนิคนับหมื่นเท่า (3)\textbf{เขียนไว้}—— พูดสั่นๆ อาจจะลืม คุณเอากระดาษมาเขียนขั้นตอนสำคัญไว้แล้วแปะไว้ด้านหลังเคสโทรศัพท์ของท่าน จะได้ผลกว่าวิธีอื่นๆ

\booksection{ขั้นตอนแรก: เปิด “โหมดใช้งานง่าย” —— โทรศัพท์จะดูสะอาดตาขึ้น}

โทรศัพท์ในปัจจุบันจะมี “โหมดใช้งานง่าย” ในตัว (บางรุ่นเรียก “โหมดดูแล” “โหมดผู้สูงอายุ” “โหมดผู้ใหญ่”) เมื่อเปิดแล้ว โทรศัพท์ทั้งเครื่องจะเปลี่ยนไปเลย:

\begin{itemize}
  \item \textbf{ฟอนต์และไอคอนใหญ่ขึ้นมาก}—— ไม่ต้องใส่แว่นสายตาอ่านอีกต่อไป
  \item \textbf{ไอคอนบนหน้าจอลดลงอย่างมาก}—— เก็บเฉพาะโทรศัพท์, ข้อความ, กล้องถ่ายรูป, อัลบั้มภาพ, WeChat และ App ที่ใช้บ่อยที่สุดไม่กี่ตัว ส่วนที่เหลือเก็บหมด
  \item \textbf{การใช้งานง่ายขึ้น}—— ฟังก์ชันต่างๆ ที่ดูซับซ้อน (หน้าจอด้านซ้ายสุด, ผู้ช่วยอัจฉริยะ, หน้าต่างลอย ฯลฯ) จะถูกปิดเป็นค่าเริ่มต้น หน้าจอจะดูสะอาดตาเป็นหลัก
  \item \textbf{เสียงดังขึ้น}—— เสียงเรียกเข้า, ระดับเสียงแจ้งเตือน, เสียงการโทร จะถูกปรับไปที่ระดับสูงสุดโดยอัตโนมัติ
\end{itemize}

\booksubsection{วิธีเปิดโหมดใช้งานง่ายของแต่ละแบรนด์ (ทำตามนี้ได้เลย):}

\begin{center}
\begin{tabularx}{\linewidth}{|X|X|X|}
\hline
\textbf{แบรนด์} & \textbf{ชื่อโหมด} & \textbf{วิธีเปิด} \tabularnewline
\hline
iPhone & การช่วยการเข้าถึง (ไม่มีโหมดใช้งานง่ายแยกต่างหาก แต่ปรับให้ใหญ่ขึ้นทีละรายการได้) & ตั้งค่า > การแสดงผลและความสว่าง > ขนาดตัวอักษร ปรับไปที่ใหญ่สุด จากนั้นเข้าไปที่ "การช่วยการเข้าถึง > การแสดงผลและขนาดตัวอักษร" เปิด "ข้อความตัวหนา" "รูปร่างปุ่ม" \tabularnewline
Huawei/Honor & โหมดใช้งานง่าย & ตั้งค่า > ระบบและการอัปเดต > โหมดใช้งานง่าย > เปิดใช้งาน \tabularnewline
Xiaomi/Redmi & โหมดขั้นต่ำ & ตั้งค่า > การตั้งค่าเพิ่มเติม > โหมดขั้นต่ำ > เปิดใช้งาน \tabularnewline
OPPO/OnePlus & โหมดใช้งานง่าย & ตั้งค่า > ฟังก์ชันเด่น > โหมดใช้งานง่าย \tabularnewline
vivo/iQOO & โหมดเดสก์ท็อปง่ายๆ & ตั้งค่า > หน้าจอหลัก, หน้าจอล็อกและวอลเปเปอร์ > โหมดเดสก์ท็อปง่ายๆ \tabularnewline
Samsung & โหมดใช้งานง่าย & ตั้งค่า > การแสดงผล > โหมดใช้งานง่าย \tabularnewline
\hline
\end{tabularx}
\end{center}

หากคุณใช้แบรนด์อื่น ให้ค้นหาคำว่า “ง่าย” หรือ “ผู้ใหญ่” ในช่องค้นหาด้านบนของ “ตั้งค่า” ในโทรศัพท์ ก็มักจะเจอ หากไม่พบ ให้ตั้งค่าทีละข้อตามแผน “ปรับใหญ่ด้วยตนเอง” ด้านล่าง ผลลัพธ์ก็จะเหมือนกัน

\booksection{หากไม่เปิดโหมดใช้งานง่าย การปรับด้วยตนเองก็ได้ —— “ตัวอักษรใหญ่+ไอคอนใหญ่” ใน 5 ขั้นตอน}

ผู้ใหญ่บางคนไม่อยากใช้โหมดใช้งานง่าย รู้สึกว่าหน้าจอดู “ล้าหลัง” เกินไป ก็ให้ปรับ 5 รายการต่อไปนี้ด้วยตนเอง ผลลัพธ์ก็ดีเช่นกัน:

\begin{enumerate}
  \item \textbf{ปรับฟอนต์ให้ใหญ่ที่สุด}:ตั้งค่า > การแสดงผล > ขนาดฟอนต์ ลากไปทางขวาสุด แล้วปรับ "ขนาดการแสดงผล" (บางรุ่นเรียก "การซูมการแสดงผล") ให้ใหญ่ที่สุดด้วย องค์ประกอบในทุกหน้าจอก็จะใหญ่ตามไปด้วย
  \item \textbf{ไอคอนใหญ่ขึ้น}:บนหน้าจอหลัก ให้ใช้สองนิ้วหนีบเข้าหากัน (เหมือนท่านซูมรูป) เข้าสู่การตั้งค่าหน้าจอหลัก > เลย์เอาต์หน้าจอหลัก (หรือขนาดไอคอน) เลือกเลย์เอาต์ที่มีช่องน้อยที่สุด (เช่น 3×5) ไอคอนก็จะใหญ่ขึ้นเอง
  \item \textbf{เปิด "ข้อความคอนทราสต์สูง"}:ทำให้ขอบของตัวอักษรคมชัดขึ้น สามารถหาสวิตช์นี้ได้ใน "การช่วยการเข้าถึง"
  \item \textbf{ปิดการแจ้งเตือนที่ไม่จำเป็น}:ตั้งค่า > การแจ้งเตือน ปิดการแจ้งเตือนของ App อื่นๆ ทั้งหมด ยกเว้นโทรศัพท์, ข้อความ และ WeChat วิธีนี้จะทำให้โทรศัพท์ไม่ดังขึ้นมาเป็นระยะ และไม่น่ารำคาญ
  \item \textbf{วางรายชื่อผู้ติดต่อที่ใช้บ่อยไว้บนหน้าจอหลัก}:กดค้างที่พื้นที่ว่างบนหน้าจอหลัก > เพิ่มวิดเจ็ต > ผู้ติดต่อ วางการโทรด่วนของคู่สมรส, ลูกๆ, แพทย์ และผู้ติดต่อฉุกเฉินไว้บนหน้าจอหลัก แค่แตะปุ่มเดียวก็โทรออกได้เลย ไม่ต้องเปิดสมุดโทรศัพท์
\end{enumerate}

\begin{tcbpractice}{ลงมือทำ: ช่วยผู้ใหญ่ "ตกแต่ง" โทรศัพท์ใหม่}
\begin{enumerate}
  \item เปิดโหมดใช้งานง่าย (ทำตามตารางแบรนด์ด้านบน)
  \item เปิดการแสดงตัวเลขแบตเตอรี่ (ตั้งค่า > แบตเตอรี่ > เปอร์เซ็นต์แบตเตอรี่) เพื่อให้ท่านรู้ตลอดว่าแบตเหลือเท่าไหร่
  \item เรียกแถบการแจ้งเตือนขึ้นมา ลบหรือซ่อนสวิตช์ลัดต่างๆ ที่ผู้ใหญ่ไม่ได้ใช้ เช่น Bluetooth, NFC, โหมดบิน
  \item วางการโทรด่วน 4 รายการบนหน้าจอหลัก: คู่สมรส, ลูกๆ, หมอประจำครอบครัว (ถ้ามี), ชุมชน/นิติบุคคล
\end{enumerate}

\end{tcbpractice}

\booksection{ป้องกันการหลอกลวง —— บทนี้โปรดอ่านไปพร้อมกับผู้ใหญ่}

ผู้สูงอายุเป็นกลุ่มเป้าหมายหลักของการหลอกลวงทางโทรศัพท์ ไม่ใช่เพราะท่านไม่ระมัดระวัง แต่เพราะคำพูดของมิจฉาชีพถูกออกแบบมาอย่างประณีต โดยมุ่งเป้าไปที่จุดอ่อนทางจิตใจของผู้สูงอายุโดยเฉพาะ ต่อไปนี้คือประเภทการหลอกลวงที่พบบ่อยที่สุดและวิธีรับมือที่ฉันได้รวบรวมไว้ ขอแนะนำให้คุณอ่านให้ผู้ใหญ่ฟัง หรือพิมพ์ออกมาแปะไว้ข้างโทรศัพท์

\booksubsection{6 วิธีหลอกลวงที่พบบ่อย (เห็นแล้วให้วางสายทันที):}

\begin{itemize}
  \item \textbf{(1) "หลานคุณเจอเรื่องแล้ว รีบโอนเงินมา"}—— มิจฉาชีพจะบอกว่า "หลานชาย/หลานสาวคุณก่อเรื่องข้างนอก โดนตำรวจจับแล้ว ต้องใช้เงินค้ำประกัน" จำไว้:\textbf{ใดๆโทรศัพท์ที่ขอให้โอนเงินเพื่อ "ปล่อยตัว" ล้วนเป็นการหลอกลวง}.ตำรวจและหน่วยงานของรัฐจะไม่โทรขอเงิน
  \item \textbf{(2) "บัตรประกันสุขภาพของคุณถูกระงับการใช้งาน"}——SMS บอกว่าบัตรประกันสุขภาพมีความผิดปกติ แล้วให้คุณคลิกลิงก์เพื่อ "เปิดใช้งาน" จำไว้:\textbf{สำนักงานประกันสุขภาพและสำนักงานประกันสังคมจะไม่ส่ง SMS ให้คุณคลิกลิงก์}.หากพบเจอสถานการณ์แบบนี้ ให้โทรไปที่ 12333 (สายด่วนทางการของสำนักงานประกันสังคม) เพื่อตรวจสอบได้เลย
  \item \textbf{(3) "การลงทุนผลตอบแทนสูง รับรองไม่ขาดทุน"}——แนะนำ "การลงทุนเพื่อผู้สูงอายุ" และ "โครงการสนับสนุนจากรัฐ" ต่างๆ ในกลุ่ม WeChat และ SMS โดยสัญญาว่าจะให้ผลตอบแทนสูง จำไว้:\textbf{ที่สัญญาผลตอบแทนต่อปีเกิน 5\% และยัง "รับประกันเงินต้น" 99.9\% คือการหลอกลวง}
  \item \textbf{(4) "ฉันคือหัวหน้า/เพื่อนเก่าของคุณ เปลี่ยนเบอร์แล้ว"}——แอบอ้างเป็นคนรู้จักมายืมเงิน จำไว้:\textbf{การโทรไปตรวจสอบคือวิธีที่ได้ผลดีที่สุด}คนที่บอกว่า "ไม่สะดวกรับสาย" แปดสิบเปอร์เซ็นต์คือมิจฉาชีพ
  \item \textbf{(5) "รับไข่ฟรี/สัมมนาผลิตภัณฑ์สุขภาพ"}——กิจกรรมฟรีแบบเจอตัวเพื่อชักชวนผู้สูงอายุฟัง "สัมมนาสุขภาพ" แล้วค่อยขายผลิตภัณฑ์สุขภาพในราคาแพงเฉียดฟ้า จำไว้:\textbf{ไม่มีของฟรีในโลก กิจกรรมฟรีทั้งหลายล้วนมีธุรกิจซ่อนอยู่}
  \item \textbf{(6) "มือถือของคุณติดไวรัสแล้ว ต้องซ่อมแบบระยะไกล"}——มิจฉาชีพจะให้คุณดาวน์โหลด App ควบคุมระยะไกล (เช่น Sunlogin, TeamViewer) แล้วเข้าควบคุมมือถือของคุณจากระยะไกลเพื่อโอนเงินไป จำไว้:\textbf{ห้ามให้ใครเข้าควบคุมมือถือของคุณจากระยะไกลเด็ดขาด!}ผู้ผลิตมือถือ ผู้ให้บริการเครือข่าย และธนาคาร จะไม่มีวันโทรมาให้คุณติดตั้งซอฟต์แวร์ควบคุมระยะไกล
\end{itemize}

\booksubsection{3 ด่านป้องกัน เมื่อตั้งค่าไว้แล้วจะช่วยบล็อกการหลอกลวงได้ถึง 90\%:}

\begin{enumerate}
  \item \textbf{ด่านที่ 1: การบล็อกในตัวมือถือ}——มือถือ Android เปิด "บล็อกสายรบกวน" (อยู่ใน App โทรศัพท์หรือ App จัดการมือถือ) ส่วน iPhone ให้เปิด "ปิดเสียงผู้โทรที่ไม่รู้จัก" (การตั้งค่า > โทรศัพท์ > ปิดเสียงผู้โทรที่ไม่รู้จัก)
  \item \textbf{ด่านที่ 2: การบล็อกโดยผู้ให้บริการเครือข่าย}——ส่ง SMS "KTFSR" ไปที่ 10086 เพื่อเปิดบริการบล็อกสายรบกวนความถี่สูงของ China Mobile ฟรี; China Telecom และ China Unicom ก็มีบริการที่คล้ายกัน แค่โทรหาฝ่ายบริการลูกค้าของผู้ให้บริการเครือข่ายเพื่อให้เจ้าหน้าที่เปิดให้ก็พอแล้ว
  \item \textbf{ด่านที่ 3: App ศูนย์ป้องกันการหลอกลวงแห่งชาติ}——ดาวน์โหลด App "National Anti-Fraud Center" (ศูนย์ป้องกันการหลอกลวงแห่งชาติ) จาก App Store แล้วลงทะเบียนและเปิดการแจ้งเตือนสายเข้า App นี้จะเชื่อมต่อกับฐานข้อมูลของตำรวจ พอเบอร์มิจฉาชีพโทรเข้ามาก็จะสามารถระบุได้ทันที
\end{enumerate}

\begin{tcbwarn}{"คำขวัญป้องกันการหลอกลวง" สำหรับผู้สูงอายุ (พิมพ์ออกมาแปะไว้ข้างโทรศัพท์)}
\begin{itemize}
  \item \textbf{โทรศัพท์แปลกหน้าขอให้โอนเงิน วางสาย!}
  \item \textbf{ลิงก์ใน SMS ให้คุณคลิก ลบทิ้ง!}
  \item \textbf{บอกว่าผลตอบแทนสูงและรับประกันเงินต้น บล็อก!}
  \item \textbf{ให้คุณติดตั้งซอฟต์แวร์ควบคุมระยะไกล วางสายทันที!}
  \item \textbf{ตำรวจของจริงจะไม่สอบสวนคดีทางโทรศัพท์และจะไม่ขอเงิน!}
  \item \textbf{เรื่องที่ยังไม่แน่ใจ ให้โทรหาลูกๆ ถามก่อน}
\end{itemize}

\end{tcbwarn}

\booksection{การช่วยเหลือจากระยะไกล — คุณอยู่ต่างจังหวัดก็ช่วยซ่อมมือถือให้พ่อแม่ได้}

บางครั้งปัญหาที่ผู้สูงอายุเจอ คุณอธิบายกันทางโทรศัพท์ไปครึ่งวัน พวกเขาก็ยังหาปุ่มที่ตรงกันไม่เจอ ในเวลาแบบนี้ฟังก์ชัน “การช่วยเหลือจากระยะไกล” ก็จะเข้ามามีประโยชน์ — คุณสามารถเห็นหน้าจอมือถือของพวกเขาบนมือถือเครื่องของคุณเอง แล้วช่วยกดให้ได้โดยตรง วิธีการของแต่ละแบรนด์จะไม่เหมือนกัน:

\begin{itemize}
  \item \textbf{Huawei/Honor}: มีแอป "Changlian" ในตัว ภายในแอปมีฟังก์ชัน "แชร์หน้าจอ" ถ้าคุณกับพ่อแม่ติดตั้งแอป Changlian ไว้ หลังจากโทรวิดีโอแล้วกด "แชร์หน้าจอ" คุณจะสามารถวงและทำเครื่องหมายบนหน้าจอของท่านได้ แล้วบอกแม่ได้ว่า "กดตรงนี้"
  \item \textbf{Xiaomi/Redmi}:ใน App วิดีโอคอลของ Xiaomi มี "การช่วยเหลือจากระยะไกล" รองรับการแชร์หน้าจอ + การคุยเสียงพร้อมกัน
  \item \textbf{OPPO/vivo}:มี "พื้นที่ครอบครัว" หรือ "ผู้พิทักษ์ระยะไกล" ในตัว ไปที่การตั้งค่าแล้วค้นหาคำว่า "ระยะไกล" ก็จะเจอ
  \item \textbf{วิธีการทั่วไป (ใช้ได้กับทุกแบรนด์)}: ดาวน์โหลดแอป "Xiangrikui Yuancheng Kongzhi" (เวอร์ชันฟรีก็ใช้ได้แล้ว) แล้วติดตั้งไว้ที่เครื่องของทั้งสองฝั่ง ลงทะเบียนด้วยบัญชีเดียวกันหรือเพิ่มเป็นเพื่อน จากนั้นให้พ่อแม่อ่านรหัสยืนยันที่ขึ้นบนหน้าจอมือถือให้ฟัง คุณก็จะสามารถดูและควบคุมหน้าจอของท่านได้จากระยะไกลแล้ว
  \item \textbf{iPhone}:ในระหว่างการโทร FaceTime รองรับ "แชร์หน้าจอ" ขณะโทรด้วย FaceTime ให้กดปุ่มแชร์หน้าจอ อีกฝ่ายก็จะสามารถเห็นหน้าจอของคุณได้ แต่เนื่องจากข้อจำกัดด้านความปลอดภัยของ iOS จะไม่สามารถควบคุมจากระยะไกลได้ สามารถ "สอนแบบสั่งการ" ได้อย่างเดียว
\end{itemize}

\booksection{การ “ปรับแต่งให้เหมาะกับผู้สูงอายุ” ของ WeChat — บอกประโยคนี้กับผู้สูงอายุก็พอแล้ว}

WeChat เป็น App ที่ผู้สูงอายุใช้บ่อยที่สุด แต่หน้าจอค่อนข้างหนาแน่นและซับซ้อนเกินไป ต่อไปนี้คือการปรับแต่งไม่กี่อย่างที่ได้ผลทันที:

\begin{itemize}
  \item \textbf{ปรับขนาดฟอนต์ของ WeChat ให้ใหญ่ขึ้น}:WeChat > ฉัน > การตั้งค่า > ทั่วไป > ขนาดฟอนต์ ลากไปที่ใหญ่ที่สุด อย่างไรก็ตาม สิ่งนี้จะมีผลเฉพาะฟอนต์ภายใน WeChat เท่านั้น จะไม่มีผลต่อฟอนต์ของระบบมือถือ
  \item \textbf{เปิด "โหมดอำนวยความสะดวก"}:WeChat > ฉัน > การตั้งค่า > โหมดอำนวยความสะดวก > เปิดใช้งาน เมื่อเปิดแล้วฟอนต์และปุ่มต่างๆ จะใหญ่ขึ้น ความคมชัดของสีจะดีขึ้น และยังสามารถกด "ฟังข้อความตัวอักษร" ได้ในคลิกเดียว (กดที่ข้อความตัวอักษรแล้วระบบจะอ่านให้ฟัง)
  \item \textbf{ช่วยตั้งชื่อให้ผู้ติดต่อสำคัญ}:เปลี่ยนชื่อของลูกๆ และญาติพี่น้องในสมุดรายชื่อให้จดจำง่าย — เช่น "ลูกชายคนโต-หมิง" "ลูกสาวคนรอง-หง" เพื่อให้พวกเขาเปิด WeChat มาแล้วรู้ทันทีว่าใครคือใคร
  \item \textbf{สอนให้พวกเขาใช้"ป้อนข้อมูลด้วยเสียง" แทนการพิมพ์}:บนแป้นพิมพ์ข้างช่องแชท ให้กดค้างปุ่มเว้นวรรค (Android) หรือแตะไอคอนไมโครโฟน (iPhone) แล้วพูดสิ่งที่ต้องการได้เลย ระบบจะแปลงเสียงเป็นข้อความเพื่อส่งให้ทันที ความแม่นยำในการจดจำเสียงตอนนี้ดีมากๆ ถือเป็นความสะดวกสบายระดับ"ปฏิวัติ" สำหรับผู้สูงอายุที่พิมพ์พินอินไม่เป็น
  \item \textbf{ปิดฟังก์ชันที่ไม่ใช้ใน"หน้าค้นพบ"}:WeChat > ฉัน > ตั้งค่า > ทั่วไป > จัดการหน้าค้นพบ ให้ปิดฟังก์ชันที่ไม่ได้ใช้ใน"ช่องคนรู้จัก""ช่องวิดีโอ""ไลฟ์สด""ช้อปปิ้ง""เกม" ทิ้งไปให้หมด เพื่อลับสิ่งรบกวน
\end{itemize}

\booksection{ทำยังไงถ้าโทรศัพท์หาย? ช่วยผู้สูงอายุเตรียม“แผนรองรับเหตุฉุกเฉิน”}

เตรียมตัวไว้ก่อนย่อมดีกว่าแก้ตอนหลัง ตอนที่โทรศัพท์ยังไม่หาย ให้ทำเรื่องเหล่านี้ให้เรียบร้อยก่อน:

\begin{enumerate}
  \item \textbf{ยืนยันว่าเปิด"ค้นหา iPhone ของฉัน" หรือ"ค้นหาอุปกรณ์" ไว้แล้ว}:iPhone อยู่ที่ การตั้งค่า > [ชื่อของคุณ] > ค้นหา > ค้นหา iPhone ของฉัน; ส่วน Android ให้ค้นหา"ค้นหาอุปกรณ์" ในการตั้งค่า ต้องแน่ใจว่าสวิตช์นี้เปิดอยู่
  \item \textbf{จดรหัสผ่านและบัญชีสำคัญไว้บนกระดาษ}:รหัสผ่านของ Apple ID / บัญชี Huawei / บัญชี Xiaomi ให้เก็บสำเนาไว้ในลิ้นชักที่บ้าน และมอบให้ลูกที่ไว้ใจได้สำเนาหนึ่ง หลังจากโทรศัพท์หาย จะต้องล็อกอินบัญชีเหล่านี้เพื่อล็อกและระบุตำแหน่งระยะไกล ถ้าลืมรหัสผ่านจะเป็นเรื่องวุ่นวายมาก
  \item \textbf{สอนประโยคเดียวให้ผู้สูงอายุ}:\textbf{"ถ้าโทรศัพท์หาย สิ่งแรกที่ต้องทำคือโทรหาครอบครัว"}ครอบครัวสามารถล็อกอิน iCloud หรือเว็บค้นหาอุปกรณ์ของ Android ทันที เพื่อตั้งค่า"โหมดสูญหาย" ระยะไกล คนที่เก็บได้จะเปิดหรือใช้งานไม่ได้เลย
\end{enumerate}

\booksection{หูฟังไร้สายและความเข้ากันได้กับเครื่องช่วยฟัง — ฟังก์ชันที่ดีแต่มักถูกมองข้าม}

ผู้สูงอายุหลายคนมีปัญหาเรื่องการได้ยิน แต่ไม่รู้ว่าโทรศัพท์สามารถทำงานร่วมกับเครื่องช่วยฟังได้ iPhone และโทรศัพท์ Android บางรุ่นรองรับเครื่องช่วยฟัง MFi (Made for iPhone) สามารถส่งเสียงจากสายตรงไปยังเครื่องช่วยฟังได้เลย ชัดเจนกว่าการถือโทรศัพท์แนบหูมาก

\begin{itemize}
  \item \textbf{วิธีเชื่อมต่อ}:iPhone ไปที่ การตั้งค่า > การช่วยการเข้าถึง > อุปกรณ์ช่วยฟัง เปิด Bluetooth แล้วทำตามคำแนะนำการจับคู่ของเครื่องช่วยฟัง
  \item \textbf{ถ้าไม่ใช้เครื่องช่วยฟัง}:ใน"การช่วยการเข้าถึง > เสียง/ภาพ" ให้เปิด"ลดสัญญาณรบกวนในการโทร" และ"เสียงโมโน" จะช่วยให้การโทรชัดเจนขึ้น
  \item \textbf{Android}:ใน"การช่วยการเข้าถึง > การช่วยฟัง" จะมีการตั้งค่าที่คล้ายกัน ตำแหน่งอาจแตกต่างกันเล็กน้อยในแต่ละแบรนด์ ให้ค้นหาด้วยคำว่า"ช่วยฟัง" ได้เลย
\end{itemize}

\booksection{ฟังก์ชันที่เกี่ยวข้องกับสุขภาพ — โทรศัพท์ไม่ใช่แค่เครื่องมือสื่อสาร}

\begin{itemize}
  \item \textbf{SOS ฉุกเฉิน}:กดปุ่มเปิด/ปิด 5 ครั้งติดต่อกัน (สำหรับโทรศัพท์ส่วนใหญ่) โทรศัพท์จะโทรหาหมายเลขฉุกเฉินโดยอัตโนมัติ และส่งพิกัดตำแหน่งให้ผู้ติดต่อฉุกเฉินที่ตั้งไว้ล่วงหน้า ช่วยผู้สูงอายุตั้งค่าไว้ก่อน — iPhone: "การตั้งค่า > SOS ฉุกเฉิน", Android: "การตั้งค่า > ความปลอดภัย > การขอความช่วยเหลือฉุกเฉิน" ก่อนสอนผู้สูงอายุใช้ ควรซ้อมลองดูสักครั้ง เผื่อตอนเกิดเหตุเขาจะได้ไม่ตื่นเต้นจนกดผิด
  \item \textbf{การตรวจจับการล้ม}:Apple Watch มีฟังก์ชันนี้ หากตรวจพบว่าคุณล้มและไม่ขยับสักพัก จะโทรหาหมายเลขฉุกเฉินและแจ้งครอบครัวโดยอัตโนมัติ ถ้ามีงบประมาณ ซื้อให้ผู้สูงอายุสักเครื่องได้ ในยามคับขันอาจช่วยชีวิตได้
  \item \textbf{การแจ้งเตือนรับประทานยา}:App "สุขภาพ" ของ iPhone มีฟังก์ชันแจ้งเตือนการรับประทานยาในตัว เพียงตั้งชื่อยา เวลา และปริมาณ เมื่อถึงเวลาจะมีการแจ้งเตือนขึ้นมา Android สามารถใช้"นาฬิกาปลุก" หรือดาวน์โหลด App ประเภท"แจ้งเตือนการกินยา" ได้
\end{itemize}

\begin{tcbsummary}{สรุปบทนี้}
\begin{itemize}
  \item เปิดโหมดง่ายๆ หรือปรับขนาดตัวอักษร/ไอคอนให้ใหญ่ขึ้นก่อน ปรับสภาพแวดล้อมให้ดีปัญหาก็จะหมดไปเกินครึ่ง
  \item กฎเหล็ก 3 ข้อในการสอนผู้สูงอายุใช้โทรศัพท์: อย่าเร่ง, ใช้การเปรียบเทียบ, จดบันทึกไว้
  \item โกง 6 รูปแบบที่พบบ่อยต้องเตือนซ้ำๆ ตั้งแนวป้องกัน 3 ชั้นไว้จะบล็อกการหลอกลวงได้ 90\%
  \item การช่วยเหลือระยะไกลช่วยให้คุณอยู่ต่างจังหวัดก็ซ่อมโทรศัพท์ให้พ่อแม่ได้
  \item การเปิด"โหมดผู้สูงอายุ" ของ WeChat + "ป้อนข้อมูลด้วยเสียง" คือ 2 ฟังก์ชันที่ใช้งานได้จริงที่สุดสำหรับผู้สูงอายุ
  \item SOS ฉุกเฉินและการตรวจจับการล้ม คือการตั้งค่าที่ช่วยชีวิตในยามคับขัน คุ้มค่าที่จะใช้เวลา 5 นาทีในการตั้งค่าให้เรียบร้อย
\end{itemize}

\end{tcbsummary}
